\subsection{Separating the archived volatility mechanisms}

The archive allows the volatility trade-off to be examined more closely.
A00 adds only the softer volatility denominator to G10; G26 switches on
the full volatility bundle while retaining G10's disagreement settings;
B02 uses the pure excess-return/volatility ratio. Across the 36 matched
cohort, breadth and allocation constructions under each entry convention,
their differences from G10 are:

\begin{minipage}{\linewidth}
\begin{center}\small\setlength{\tabcolsep}{4pt}\renewcommand{\arraystretch}{1.18}
\begin{longtable}{@{}lrrrr@{}}\toprule
Rule & Anchor R pp & Anchor DD gain & Available R pp & Available DD gain\\\midrule\endhead
A00 & -0.49 & +1.21 & +0.04 & +0.53\\
G26 & -4.51 & +3.71 & -1.80 & +1.61\\
B02 & -6.17 & +4.99 & -1.82 & +3.14\\
\bottomrule\end{longtable}\end{center}
\end{minipage}\par

A00 preserves more upside while providing a modest drawdown improvement.
B02 makes the stronger risk-objective trade: its average drawdown gain is
4.99 percentage points at anchor entry and 3.14 points at availability,
compared with 3.71 and 1.61 points for the bundled G26 rule. The profiles
also differ in design: A00's denominator rises only to 1.25, versus 2.0
inside the full bundle, and B02 omits the ensemble-disagreement penalties.
This comparison explains the separate roles of the simple return-focused
and volatility-adjusted screens.
